In this section, we prove Theorem~\ref{thm:virtual2system}, which we
repeat here for convenience.

\medskip
\noindent
	\textbf{Theorem~\ref{thm:virtual2system}} (\textbf{Groups expressed by linear systems of bounded size}).
\hfill\\
	\emph{Let $L$ be the group spanned by a finite set $X\subseteq
	\setZ^d$. There exists a~$d$-dim linear system $S$ such that
	\begin{center}
	$\sem{S}=L$ and $\max\{\nnorm{S},\mlcm{S}\}\leq d\,!\cdot \norm{}{X}^d $.
	\end{center}
		}
\medskip

\noindent
Recall that the group $L$ spanned by a finite set $X\subseteq \setZ^d$ is 
the set of integer linear combinations of elements of $X$. The
elements of $X$ can be viewed as the columns of 
an integer $d\times k$ matrix $C$, where $k=|X|$; in this case,
$\nnorm{X}=\nnorm{C}$. 

To prove the theorem, we consider an integer $d\times k$ matrix $C$
and 
denote by $L_C$ the
 group (also called a \emph{lattice}) spanned by the columns
of $C$; that is,
\begin{equation*}
L_C=\{x\in\setZ^d\mid \exists y\in \setZ^k: Cy=x\}.
\end{equation*}
We aim to construct a linear system $S$ such that $\sem{S}=L_C$ and
\begin{equation}\label{eq:boundSLC}
	\max\{\nnorm{S},\mlcm{S}\}\leq d\,!\cdot \nnorm{C}^d.
\end{equation}	
\paragraph{Case $\bm{\text{rank}(C)=d}$.}
We first consider the case where $\rank{C} = d$; that is, we assume that the rows of $C$ are linearly independent. From~\cite[Chapter 4]{schrijver1986theory}, we recall that there exists a square $d \times d$ non-singular integer matrix $H$ such that $L_H = L_C$. We can take $H$ to be the first $d$ columns of the Hermite normal form of $C$ (the remaining $k-d$ columns being zero vectors). It follows that 
\[
    L_C = L_H= \{x\in\setZ^d\mid \exists y\in \setZ^d: Hy=x\}=\{x \in \setZ^d \mid H^{-1}x \in \setZ^d\};
\]
note that $H^{-1}x \in \setZ^d$ implies that $x = H(H^{-1}x) \in \setZ^d$.

Although the entries of $H^{-1}$ may be non-integers, the matrix
$\det(H) \cdot H^{-1}$ is an integer matrix, since $H^{-1} =
\frac{1}{\det(H)}\cdot \adj(H)$, where $\adj(H)$ is the adjugate of
$H$ (and thus an integer matrix). In particular, $H^{-1}(\det(H) \cdot
z)$ is an integer vector for any $z \in \setZ^d$. This implies that
for any $x \in \setZ^d$, 
\[
	\text{the vector } H^{-1}x \text{ is an integer vector if and
	only if } H^{-1}(x \bmod \ell) \text{ is an integer vector,}
\]
where $\ell = |\det(H)|$ and the modulo operation is applied component-wise.


Defining $B = \{b \in \{0, 1, \dots, \ell{-}1\}^d \mid H^{-1}b \in \mathbb{Z}^d\}$, we obtain $L_C = \sem{S}$, where 
\[
    S = \bigvee_{b \in B} \left( \bigwedge_{i=1}^d \bvar{x}_i \equiv
    b(i) \!\!\!\! \pmod \ell \right);
\]
here $\bvar{x} = (\bvar{x}_1,\bvar{x}_2, \dots, \bvar{x}_d)$ is the
variable vector of $S$ ranging over $\setZ^d$.

To ensure the bound~(\ref{eq:boundSLC}),
it suffices to show that $|\det(H)|\leq d\,!\cdot\nnorm{C}^d$.
Consider any $d\times d$ matrix $M$ composed of $d$
linearly independent columns of $C$; hence $L_M$ is a subgroup of
$L_C=L_H$.
From~\cite[Chapter 4]{schrijver1986theory}, it follows that $\det(H)$
divides $\det(M)$, which implies that $|\det(H)|\leq |\det(M)|$.
Since $|\det(M)|\leq d\,!\cdot\nnorm{M}^d$ and
$\nnorm{M}\leq \nnorm{C}$, we have completed the case for
$\rank{C}=d$.

\paragraph{Case $\bm{\text{rank}(C)=r<d}$.}
We now consider a $d \times k$ integer matrix $C$ with $\rank{C} = r < d$. Without loss of generality, we assume (possibly after reordering the rows of $C$) that $C$ is of the form
\[
	C = \begin{bmatrix} \bar{C} \\ C' \end{bmatrix},
\]
where $\bar{C}$ is an $r \times k$ matrix with $\rank{\bar{C}} = r$
(that is, $\bar{C}$ is a row basis of $C$), and each of the $d-r$ rows
of $C'$ is a rational linear combination of the $r$ rows of $\bar{C}$.
Hence, there exists a unique rational $(d-r) \times r$ matrix $T$ such
that $C' = T\bar{C}$.

We now examine the transformation of $\bar{C}$ into its Hermite normal
form in more detail. This is achieved by a sequence of elementary
column operations on $\bar{C}$, comprising: 
\begin{enumerate*}[(i)]
    \item multiplying a column by $-1$, 
    \item swapping two columns, and 
    \item adding an integer multiple of one column to another.
\end{enumerate*}
These operations correspond to multiplying $\bar{C}$ from the right by a unimodular matrix $U$ (an integer $k \times k$ matrix with $|\det(U)| = 1$). Multiplying $C$ by $U$, we obtain
\[
	CU = \begin{bmatrix} \bar{C} \\ C' \end{bmatrix} U = 
		\begin{bmatrix} \bar{H} & 0 \\ H' & 0 \end{bmatrix} = [H \mid 0],
\]
where $\bar{H}$ is a non-singular $r \times r$ integer matrix such
that $L_{\bar{H}}
= L_{\bar{C}}$, and $H' = T\bar{H}$. The latter holds because $C' = T\bar{C}$ and elementary column operations preserve linear relations between rows.

Since elementary column operations preserve the spanned group, we have
\[
    L_{C} = L_{H} = \{x \in \setZ^d \mid \exists y \in
    \setZ^r : \bar{H}y = x\restr{[1,r]} \land H'y = x\restr{[r+1,d]}\}.
\]
By rewriting the conjunction as $y = \bar{H}^{-1}x\restr{[1,r]}\land
H'\bar{H}^{-1}x\restr{[1,r]} = x\restr{[r+1,d]}$, and noting that
$H'\bar{H}^{-1} = T$, we obtain
\begin{equation}\label{eq:LbarC}
	L_{C}=\{x\in\setZ^d\mid x\restr{[1,r]}\in L_{\bar{C}}\land
x\restr{[r+1,d]}=T x\restr{[1,r]}\}.
\end{equation}

% Fix Indentation
\noindent
We note that the rational $(d{-}r)\times r$ matrix $T$ is determined
by any submatrix $M$ consisting of 
$r$ linearly independent columns of $C$.
Letting
$\bar{M}=M\restr{[1,r]}$
and $M'=M\restr{[r+1,d]}$, we have $M'=T\bar{M}$ (as $C'=T\bar{C}$), which
implies that the $r$ rows of $\bar{M}$ are linearly independent;
hence, $\bar{M}$
is non-singular.
 Using Cramer's rule,
we deduce that all entries of $T$ are of the form
$\frac{\det(M'')}{\det(\bar{M})}$, with the same denominator
$\det(\bar{M})$ and
varying integer numerators. The absolute values of the denominator and all numerators are
bounded by $r\,!\cdot \nnorm{M}^r$, where $r<d$ and $\nnorm{M}\leq \nnorm{C}$.

The condition $x\restr{[r+1,d]}=T x\restr{[1,r]}$ can be expressed
as a conjunction of homogeneous linear constraints for $i\in[r{+}1,d]$ of the form 
\[
	\det(\bar{M}) \cdot \bvar{x}_i - \sum_{j=1}^r \det(M''_{ij}) \cdot
    \bvar{x}_j = 0,
\]
where the absolute values of all coefficients are bounded by
$r\,! \cdot \nnorm{C}^r$. 
Since $r < d$, the group $L_{C}$ in~\eqref{eq:LbarC} can indeed
be expressed by a linear system $S$ satisfying~\eqref{eq:boundSLC}.
Here, we apply the results from the case $\rank{C} = d$ to the
submatrix $\bar{C}$ with $\rank{\bar{C}} = r$ (noting that
$\nnorm{\bar{C}} \leq \nnorm{C}$). This completes the proof of Theorem~\ref{thm:virtual2system}.
